\documentclass{article}
\usepackage{mathnotes}

\title{Some Discrete Probability Distributions}
\description{Analysis of binomial, hypergeometric, Poisson, geometric, and negative binomial distributions with their probability mass functions and parameters.}

\begin{document}

\section{Some Discrete Probability Distributions}

\subsection{Binomial Distributions}

\begin{definition}[Bernoulli Process]
A \textbf{Bernoulli Process} must possess the following properties:

\begin{enumerate}
\item The experiment consists of repeated trials.
\item Each trial results in a boolean outcome, which can be considered true or false, success or failure, etc.
\item The probability of success, $p$, remains constant across trials.
\item Repeated trials are independent.
\end{enumerate}
\end{definition}

\begin{definition}[Binomial Distribution]
A \textbf{binomial distribution} aggregates the outcomes of a \@{Bernoulli-process} across multiple trials to tell you how likely it is to achieve a certain number of successes. The number $X$ of successes in $n$ Bernoulli trials is called a \textbf{binomial random variable,} and its probability distribution is what we call the binomial distribution.

If we let $n$ be the number of trials, $x$ be the number of successes, and $p$ be the probability of success on a given trial, and $q = 1 - p$ be the probability of failure on a given trial, then the probability distribution of the random variable $X$ is

\[ b(x, n, p) = {n \choose x}p^x q^{n-x}, \quad x = 0, 1, 2, \dots, n. \]

To break this down some, there are ${n \choose x}$ ways to have $x$ successes out of $n$ trials (different orderings). For each, there are $x$ independent events that occur with a probability of $p$ and $n - x$ independent events that occur with a probability of $q = 1 - p.$

The mean and \@{variance} of the binomial distribution $b(x, n, p)$ are

\[ \mu = np, ~ \text{and} ~ \sigma^2 = npq. \]

We can use a summation over the formula above to find the probability of there being between $a$ and $b$ successes:

\[ \sum_{x=a}^b{\binom{n}{x} p^x q^{n-x}}. \]
\end{definition}

\subsection{Hypergeometric Distributions}

\begin{definition}[Hypergeometric Distribution]
The Hypergeometric distribution is similar to the \@{binomial-distribution}, but is performed without replacement. So, if success is drawing an ace from a deck of cards, in a Binomial situation the card drawn each trial would be put back into the deck; in the Hypergeometric it would not be. Thus, the trials for a Hypergeometric distribution are not independent.

A \textbf{Hypergeometric experiment} has the following properties:

\begin{enumerate}
\item A random sample of size $n$ is selected without replacement from $N$ items.
\item Of the $N$ items, $k$ may be classified as successes and $N - k$ are classified as failures.
\end{enumerate}

The number $X$ of successes of a Hypergeometric experiment is called a \textbf{Hypergeometric random variable,} and its probability distribution is called the \textbf{Hypergeometric distribution.}

The probability distribution of $X$, the probability of $x$ successes in $n$ draws, without replacement, from a finite population of size $N$ that contains exactly $k$ successful items and $N - k$ failure items, is

\[ h(x, N, n, k) = \frac{\binom{k}{x}\binom{N-k}{n-x}}{\binom{N}{n}}, \quad \max{(0, n - (N -k))} \leq x \leq \min{(n, k)}. \]

The mean and \@{variance} of the Hypergeometric distribution $h(x, N, n, k)$ are

\[ \mu = \frac{nk}{N} ~ \text{and} ~ \sigma^2 = \frac{N - n}{N - 1} \cdot n \cdot \frac{k}{N} \left ( 1 - \frac{k}{N} \right ). \]

It's worth noting that when $n$ is small compared to $N$, the lack of replacement in a hypergeometric process doesn't cause much impact to the distribution, and in these cases, the hypergeometric distribution is similar to the binomial distribution. In fact, for large values of $N$ and small values of $n$, the binomial distribution approximates the hypergeometric distribution.

To do this approximation, just use the binomial PMF with $p = k/N.$
\end{definition}

\subsection{Poisson Distribution}

\begin{definition}[Poisson Distribution]
Experiments that give numerical values of a random variable $X,$ the number of outcomes during a given time interval or in a specified region, are called \textbf{Poisson experiments.} For example, the number of phone calls per hour an office receives or the number of field mice per acre in a pasture. Note that while binomial and hypergeometric experiments dealt with discrete sequence of events (and give probabilities for discrete outcomes), Poisson experiments deal with continuous domains (and also give probabilities for discrete outcomes).

A \textbf{Poisson process} possesses the following properties

\begin{enumerate}
\item The number of outcomes occurring in one time/space region is independent of the number of outcomes occurring in any other disjoint time/space region, i.e., the Poisson process has no memory.
\item The probability that a single outcome will occur within a small region is proportional to the size of the region and does not depend on the number of outcomes occurring outside the region, or the relative position of the region.
\item In a very small region, the probability of more than one event occurring is negligible.
\end{enumerate}

The probability distribution of the Poisson random variable $X,$ representing the number of outcomes occurring in a given time interval or specified region denoted by $t,$ is

\[ p(x, \lambda t) = \frac{e^{-\lambda t}(\lambda t)^x}{x!}. \]

Both the mean and \@{variance} of the Poisson distribution $p(x, \lambda t)$ are $\lambda t.$
\end{definition}

Given a binomial distribution, if $n$ is large and $p$ is close to 0, the Poisson distribution approximates binomial probabilities.

\begin{theorem}[Poisson Approximation to the Binomial]\label{poisson-approximation-binomial}
Let $X$ be a \@[binomial]{binomial-distribution} random variable with probability distribution $b(x, n, p).$ When $n \to \infty,$ $p \to 0,$ and $\lim_{n \to \infty}{np} = \mu$ remains constant,

\[ \lim_{n \to \infty}{b(x, n, p)} = p(x, \mu). \]
\end{theorem}

So, to do this approximation, let $\lambda t = \mu = np.$

Note: some texts just use $\lambda$ in place of $\lambda t.$

\subsection{Geometric Distribution}

\begin{definition}[Geometric Distribution]
A \textbf{geometric random variable} is a discrete random variable that models the number of independent Bernoulli trials needed to achieve the first success. The probability of success is constant (denoted by $p$) and the trials are independent. The probability mass function is

\[ P(X = k) = (1 - p)^{k-1} p, \]

for $k = 1, 2, 3 \dots.$ $k$ represents the number of independent trials it takes to achieve the first success, and $p$ is the probability of any given trial being successful.

The mean is $\frac{1}{p}$ and the \@{variance} is $\frac{1 - p}{p^2}.$
\end{definition}

\subsection{Negative Binomial Distribution}

\begin{definition}[Negative Binomial Distribution]
The \textbf{negative binomial random variable} generalizes the \@[geometric]{geometric-distribution} random variable by giving the probability that it takes $n$ independent trials to get $r$ successes. Its pmf is

\[ P(X = n) = \binom{n - 1}{r - 1} p^r (1 - p)^{n - r}. \]
\end{definition}

\end{document}
